\begin{description}
\item[(a)] At radius $r_0$, the angular velocity is
  \[ \Omega_0 = \sqrt{\frac{GM}{r_0^3}} \]
  \hfill(0.5 points)

  The gravitational potential energy of the particle is
  \[ \Phi_{\mathrm{grav}}(r, z) = -\frac{GMm}{\sqrt{r^2 + z^2}} \]
  \hfill(0.5 points)

  The centrifugal potential energy of the particle is
  \[ \Phi_{\mathrm{cent}}(r) = -\frac{1}{2}m\Omega_0^2 r^2 \]
  \hfill(0.5 points)

  The effective potential energy of the particle is
  \[ \Phi_{\mathrm{eff}}(r, z)
     = \Phi_{\mathrm{grav}}(r, z) + \Phi_{\mathrm{cent}}(r)
     = -\frac{GMm}{\sqrt{r^2 + z^2}} - \frac{1}{2}m\Omega_0^2 r^2 \]
  \hfill(0.5 points)

\item[(b)] Using the variable $x$, the coordinates of the particle in the
  vicinity of $r_0$ can be written as
  \[ r = r_0(1 + x), \qquad z \approx r_0 x\tan\theta_0 \]
  \hfill(0.5 points)

  The gravitational potential energy expressed in terms of the variable $x$
  is
  \[ \Phi_{\mathrm{grav}}(x)
     = -\frac{GMm}{r_0}\left[1 + 2x + x^2(1 + \tan^2\theta_0)\right]^{-\frac{1}{2}} \]
  \hfill(0.5 points)

  Defining $u = 2x + x^2(1 + \tan^2\theta_0)$ and using the Taylor expansion
  $(1 + u)^{-1/2} \approx 1 - \frac{u}{2} + \frac{3}{8}u^2$ while retaining
  only terms up to second order in $x$, yields
  \begin{align*}
    \Phi_{\mathrm{grav}}(x)
      &\approx -m\Omega_0^2 r_0^2\left[1 - \frac{1}{2}
         \left(2x + x^2 + x^2\tan^2\theta_0\right)
         + \frac{3}{8}\left(4x^2\right)\right] \\
      &\approx -m\Omega_0^2 r_0^2\left[1 - x
         + x^2\left(1 - \frac{\tan^2\theta_0}{2}\right)\right]
  \end{align*}
  \hfill(1.5 points)

  The centrifugal potential energy expressed in terms of the variable $x$ is
  \[ \Phi_{\mathrm{cent}}(x) = -\frac{1}{2}m\Omega_0^2 r_0^2(1 + x)^2
     = -m\Omega_0^2 r_0^2\left(\frac{1}{2} + x + \frac{1}{2}x^2\right) \]
  \hfill(0.5 points)

  The effective potential energy expressed in terms of the variable $x$ is
  \[ \Phi_{\mathrm{eff}}(x) = -m\Omega_0^2 r_0^2\left[\frac{3}{2}
     + x^2\left(\frac{3}{2} - \frac{\tan^2\theta_0}{2}\right)\right] \]
  \hfill(1 point)

  The effective potential energy has a parabolic form and reaches an
  extremum at $x = 0$. For the particle to move away from the disk along the
  magnetic field line, this extremum must be a maximum of effective
  potential energy. This condition leads to
  \[ 3 - \tan^2\theta_0 > 0 \Rightarrow \theta_0 < 60^\circ \]
  % sic: the degree sign is printed as a superscript zero (60^0), here and in the note below
  \hfill(1 point)

  \emph{Note: Students can find $\theta_0 < 60^\circ$ from
  $F_{\mathrm{cent}} > F_{\mathrm{grav}}$, where $F_{\mathrm{cent}}$ and
  $F_{\mathrm{grav}}$ are centrifugal and gravitational forces,
  respectively.}

\item[(c)] At radius $r = r_{\mathrm{A}}$, the particles still rotate at
  angular velocity $\Omega_0$.

  The rate of angular momentum removal by the stellar wind is
  \[ \Delta L_{\mathrm{w}} = \Delta M_{\mathrm{w}}\Omega_0 r_{\mathrm{A}}^2 \]
  \hfill(1 point)

  The angular momentum releasing rate by accretion is
  \[ \Delta L_{\mathrm{acc}} = \Delta M_{\mathrm{acc}}\Omega_0 r_0^2 \]
  \hfill(1 point)

  Consider all the angular momentum released by accretion is removed through
  the stellar wind, we have
  \[ \Delta L_{\mathrm{w}} = \Delta L_{\mathrm{acc}}
     \Rightarrow \Delta M_{\mathrm{w}}\Omega_0 r_{\mathrm{A}}^2
       = \Delta M_{\mathrm{acc}}\Omega_0 r_0^2
     \Rightarrow \frac{\Delta M_{\mathrm{acc}}}{\Delta M_{\mathrm{w}}}
       = \frac{r_{\mathrm{A}}^2}{r_0^2} \]
  \hfill(1 point)
\end{description}
